% TOP_PROBLEM: {"id": 20000677, "title": "Correlation cancellation for multiplicative functions", "category_id": 1, "category": {"id": 1, "name": "number_theory", "display_name": "Number Theory", "description": "Properties of integers, prime numbers, Diophantine equations.", "slug": "number-theory", "order_index": 1, "created_at": "2026-07-31T15:26:25.670Z"}, "status": "solved"}
% FIRST_POSTED: null
% ENTRY_KIND: statement_only
% !TeX program = lualatex
\documentclass[11pt]{article}
\usepackage[margin=25mm]{geometry}
\usepackage{fontspec}
\setmainfont{Latin Modern Roman}
\usepackage{amsmath,amssymb,mathtools}
\usepackage{unicode-math}
\setmathfont{Latin Modern Math}
\usepackage{xurl,hyperref}
\hypersetup{colorlinks=true,urlcolor=blue}
\setcounter{secnumdepth}{0}
\setlength{\parindent}{0pt}
\setlength{\parskip}{5pt}
\setlength{\emergencystretch}{3em}
\newcommand{\R}{\mathbb R}
\newcommand{\N}{\mathbb N}
\newcommand{\Z}{\mathbb Z}
\newcommand{\Q}{\mathbb Q}
\newcommand{\C}{\mathbb C}
\newcommand{\D}{\mathbb D}
\newcommand{\F}{\mathbb F}
\newcommand{\sref}[2]{\hyperlink{source:#1:#2}{[S#2]}}
\newcommand{\cataloguescope}[1]{#1}
\begin{document}
% UnsolvedMath ID 20000677; source catalogue code AIM-ANALYTIC_NUMBER_THEORY-0041
\setcounter{section}{20000676}
\section{Correlation cancellation for multiplicative functions}\label{20000677}
{\small UnsolvedMath ID 20000677 · Number Theory}
\subsection{Definitions and mathematical statement}

\paragraph{Shared notation.}$\mathbb N$, $\mathbb Z$, $\mathbb Q$, $\mathbb R$, and $\mathbb C$ denote the natural numbers, integers, rationals, reals, and complex numbers, respectively. All additional parameters and hypotheses are specified in the question.


\paragraph{Statement recorded in the supplied source.}
Suppose that for all fixed $\chi$ mod $q$ and $t \in \R$, $f:\N \to \C$ is multiplicative with
\begin{equation*}
 \sum_{p} \frac{1 - \operatorname{Re} f(p) \overline{\chi(p) p^{it}}}{p} = \infty.
\end{equation*}
Does there exists a subsequence $X_k \to \infty$ such that
\begin{equation*}
 \sum_{n \leq X_j} f(n) \overline{f(n+1)} = o(X_j)?
\end{equation*}
Alternatively, is
\begin{equation*}
 \sum_{n\leq X} f(n) f(n+h_1) \cdots f(n+h_k) = o(X)
\end{equation*}
for distinct $h_1, \dots, h_k$?
\par\smallskip\textit{Formulation references:} \sref{20000677}{1}, \sref{20000677}{2}, \sref{20000677}{3}.
\subsection{Short English statement}
Suppose that for all fixed $\chi$ mod $q$ and $t \in \R$, $f:\N \to \C$ is multiplicative with
\begin{equation*}
 \sum_{p} \frac{1 - \operatorname{Re} f(p) \overline{\chi(p) p^{it}}}{p} = \infty.
\end{equation*}
Does there exists a subsequence $X_k \to \infty$ such that
\begin{equation*}
 \sum_{n \leq X_j} f(n) \overline{f(n+1)} = o(X_j)?
\end{equation*}
Alternatively, is
\begin{equation*}
 \sum_{n\leq X} f(n) f(n+h_1) \cdots f(n+h_k) = o(X)
\end{equation*}
for distinct $h_1, \dots, h_k$?
\subsection{Sources}
\begin{itemize}
\item[\hypertarget{source:20000677:1}{S1}] ULAM.AI and the UnsolvedMath Contributors, UnsolvedMath: A Curated Collection of Open Mathematics Problems, record 20000677 (AIM-ANALYTIC\_NUMBER\_THEORY-0041). Catalogue attribution under CC BY 4.0.
\url{https://huggingface.co/datasets/ulamai/UnsolvedMath}.
\item[\hypertarget{source:20000677:2}{S2}] Original problem formulation and mathematical context.
\url{http://aimpl.org/sarnakconjecture/7/}.
\item[\hypertarget{source:20000677:3}{S3}] Original problem formulation and mathematical context.
\url{https://arxiv.org/abs/2304.05344}.
\end{itemize}
\label{end:20000677}
\end{document}
